%% =====================================================================
%%  B5-T-25-03 -- what B5 contributes to "The Anchor"
%%  -------------------------------------------------------------------
%%  LAYER A: APPEND-ONLY. Written once, then left alone. Material found
%%  only here sits in the `unique` slot and is protected by rule R10.
%% =====================================================================
%% @kind: source
%% @id: B5-T-25-03
%% @book: B5
%% @topic: T-25-03
%% @locator: ch. 6, pp. 152--187
%% @status: distilled
%% @unique: yes
%% @integrated: yes
%% @updated: 2026-09-19

\dossier{B5}{T-25-03}{\bookshort{B5}}{ch. 6, pp. 152--187}

\begin{distilled}
  \item The chapter is built on the Kahneman--Tversky line: people evaluate outcomes from a reference point, and losses hurt roughly twice as much as equivalent gains please (loss aversion) --- so a counterpart shown a plausible loss moves to avoid it faster than to secure a gain.
  \item \emph{Anchor their emotions}: before any number, name their fears (an accusation audit), anchoring their emotional state in preparation for a loss so they will jump at avoiding it.
  \item On numbers: going first rarely helps, but a \emph{range} can be offered without conceding --- recall a similar deal to establish the best believable ballpark, aimed at your real target.
  \item The word \emph{fair} is analysed as the F-word of negotiation, with three uses: as pressure (we just want a fair deal, stated before you know what they want), as accusation (that's not fair, an attack dressed as principle), and as the legitimate positive (proposing the deal be judged fair by outside standards).
  \item Compromise is attacked head-on: splitting the difference is called an easy outcome, not a right one, with the ransom case showing a midpoint that is bizarrely bad for both sides.
  \item The advice set also includes odd-number anchors (precise figures read as calculated and real) and letting a stalemate sit rather than smoothing it.
\end{distilled}

\begin{unique}
  \item The emotional anchor preceding the numeric one: the counterpart's fear is priced before the price is, which is why the number later lands as relief.
  \item The fair taxonomy as a live recognition tool: all three uses are moves on the frame, and each has a distinct counter.
  \item The range technique as the answer to who-anchors-first: a range sets the reference while appearing to respond.
\end{unique}
